De forma a distribuir o ESPEON e executar os testes em sala de aula remota, faz-se necessário o \textit{deploy} da aplicação em ambiente produtivo, acessível pela \textit{internet}. Para fazer isso foram escolhidos dois serviços de hospedagem gratuitos, Render e AtlasDB. O Render contempla o \textit{backend} ESPEON e o banco PostgreSQL, conforme indicado na Figura \ref{fig:render_stacks}

\begin{figure}[H]
    \centering
    \begin{minipage}{0.7\textwidth}
        \centering
        \includegraphics[width=.8\textwidth]{espeon/suporte.render.png}
        \caption{Suporte Render. \textbf{fonte: \href{https://render.com/docs/faq}}{Render.com}} 
        \label{fig:render_stacks}
    \end{minipage}\hfill
\end{figure}

Ainda que seja possível implantar o banco MongoDB, não é um processo gerenciado nativamente pelo serviço \textit{web}, então --- para evitar enventuais complicações --- optou-se pela hospedagem AtlasDB, com suporte nativo para o banco NoSQL que pretendemos implantar. No entanto, Atlas possui controle de acesso por lista de IPs (Figura~\ref{fig:atlas_whitelist} ), que limita conexões que não sejam da sua máquina, inicalmente. 

\begin{figure}[H]
    \centering
    \begin{minipage}{0.9\textwidth}
        \centering
        \includegraphics[width=.8\textwidth]{espeon/atlas.whitelist.png}
        \caption{Habilitando acesso da origem Render} 
        \label{fig:atlas_whitelist}
    \end{minipage}\hfill
\end{figure}

Então, para possibilitar a comunicação entre Render e Atlas, que ocorre quando o ESPEON API tenta gravar um \textit{log} de atividade do aluno, deve-se incluir os IPs de saída da instância Render na \textit{whitelist} do Atlas.

